\documentclass{article}

\usepackage{amsmath,amssymb}
\usepackage{xurl}
\usepackage[hidelinks]{hyperref}
\setlength{\emergencystretch}{2em}

\input{command}

\title{The Minkowski Canonical-Pair Dependency in Four Real Dimensions}
\author{}
\date{2026-08-24}

\begin{document}
\maketitle
\gapnote{open-gap}{wip}

\section{The required assertion}

Define the Minkowski metric
\begin{align}
    M
        &=\operatorname{diag}(1,-1,-1,-1).
    \label{eq:qiclean_minkowski_canonical_pair_dependency_minkowski_metric}
\end{align}
Let $C\in M_4(\mathbb R)$ be self-adjoint with respect to $M$, meaning that
\begin{align}
    C^{\mathsf T}M
        &=MC.
    \label{eq:qiclean_minkowski_canonical_pair_dependency_m_selfadjoint}
\end{align}
The proof of Theorem~3 of Verstraete--Dehaene--De Moor invokes
Gohberg--Lancaster--Rodman, Part~I, Chapter~5, Section~5.1, Theorem~5.3.
In the convention required here, that result produces an invertible real
matrix $X$, a real Jordan matrix $J$, and a canonical metric $N_J$ satisfying
\begin{align}
    C
        &=X^{-1}JX,
    \label{eq:qiclean_minkowski_canonical_pair_dependency_similarity}\\
    XMX^{\mathsf T}
        &=N_J.
    \label{eq:qiclean_minkowski_canonical_pair_dependency_metric_congruence}
\end{align}

For a real-eigenvalue Jordan block of size $n$, the corresponding metric block
is $\varepsilon P_n$, where $\varepsilon\in\{1,-1\}$ and $P_n$ is the reversal
permutation matrix. A block for a non-real conjugate pair has real dimension
$2m$ and an unsigned reversal metric; it has no sign characteristic.

The convention in
(\ref{eq:qiclean_minkowski_canonical_pair_dependency_similarity})--%
(\ref{eq:qiclean_minkowski_canonical_pair_dependency_metric_congruence})
is obtained by applying the cited canonical-pair theorem to
$C^{\mathsf T}$, rather than directly to $C$. Indeed,
(\ref{eq:qiclean_minkowski_canonical_pair_dependency_m_selfadjoint}) and
$M^2=I$ imply
\begin{align}
    CM
        &=MC^{\mathsf T}.
    \label{eq:qiclean_minkowski_canonical_pair_dependency_transpose_selfadjoint}
\end{align}
The cited theorem then gives an invertible real matrix $S$, a real Jordan
matrix $J$, and a canonical signed metric $P_{\varepsilon,J}$ such that
\begin{align}
    C^{\mathsf T}
        &=S^{-1}JS,
    \label{eq:qiclean_minkowski_canonical_pair_dependency_glr_similarity}\\
    M
        &=S^{\mathsf T}P_{\varepsilon,J}S.
    \label{eq:qiclean_minkowski_canonical_pair_dependency_glr_congruence}
\end{align}
Set $X=S^{-\mathsf T}$. Transposing
(\ref{eq:qiclean_minkowski_canonical_pair_dependency_glr_similarity}) and
using
(\ref{eq:qiclean_minkowski_canonical_pair_dependency_glr_congruence}) gives
\begin{align}
    C
        &=S^{\mathsf T}J^{\mathsf T}S^{-\mathsf T},
    \label{eq:qiclean_minkowski_canonical_pair_dependency_transposed_similarity}\\
    C
        &=X^{-1}J^{\mathsf T}X,
    \label{eq:qiclean_minkowski_canonical_pair_dependency_x_similarity}\\
    XMX^{\mathsf T}
        &=S^{-\mathsf T}MS^{-1},
    \label{eq:qiclean_minkowski_canonical_pair_dependency_x_congruence_step}\\
    XMX^{\mathsf T}
        &=P_{\varepsilon,J}.
    \label{eq:qiclean_minkowski_canonical_pair_dependency_x_congruence}
\end{align}

It remains to restore the chosen upper-Jordan orientation. For every real
Jordan block $J_k(\lambda)$, reversal satisfies
\begin{align}
    P_kJ_k(\lambda)^{\mathsf T}P_k
        &=J_k(\lambda),
    \label{eq:qiclean_minkowski_canonical_pair_dependency_real_block_reversal}\\
    P_k^2
        &=I.
    \label{eq:qiclean_minkowski_canonical_pair_dependency_reversal_involution}
\end{align}
Simultaneously reversing the chain and realification coordinates gives the
corresponding identity for each non-real conjugate-pair block. Let $R$ be the
direct sum of these blockwise reversals. Replacing $J^{\mathsf T}$ by
$RJ^{\mathsf T}R$ and $X$ by $RX$ preserves
(\ref{eq:qiclean_minkowski_canonical_pair_dependency_x_similarity}). It also
preserves each metric block:
\begin{align}
    R(\varepsilon P_k)R^{\mathsf T}
        &=\varepsilon P_k
    \label{eq:qiclean_minkowski_canonical_pair_dependency_signed_metric_reversal}
\end{align}
for a real block, with the same identity for an unsigned conjugate-pair block.
Thus $N_J:=P_{\varepsilon,J}$ has the required orientation and sign
convention. These inverse, transpose, and reversal operations are essential.
An ordinary Euclidean singular value decomposition does not prove
(\ref{eq:qiclean_minkowski_canonical_pair_dependency_similarity})--%
(\ref{eq:qiclean_minkowski_canonical_pair_dependency_metric_congruence}).

\section{The four-dimensional inertia restriction}

The reversal form $P_n$ has inertia
\begin{align}
    \operatorname{inertia}(P_n)
        &=\left(\left\lceil\frac{n}{2}\right\rceil,
            \left\lfloor\frac{n}{2}\right\rfloor\right).
    \label{eq:qiclean_minkowski_canonical_pair_dependency_reversal_inertia}
\end{align}
Multiplication by a negative sign interchanges the two indices. An unsigned
conjugate-pair block of real dimension $2m$ has inertia
\begin{align}
    \operatorname{inertia}(P_{2m})
        &=(m,m).
    \label{eq:qiclean_minkowski_canonical_pair_dependency_complex_pair_inertia}
\end{align}
Total inertia $(1,3)$ therefore excludes a real Jordan block of size four, a
conjugate-pair chain of real dimension four, and two blocks of dimension two.
The remaining dimension patterns are exactly those used in the proof of
Theorem~3 under review:
\begin{enumerate}
    \item four real one-dimensional blocks;
    \item one two-dimensional non-real conjugate-pair block and two real
          one-dimensional blocks;
    \item one real Jordan block of size two and two real one-dimensional
          blocks;
    \item one real Jordan block of size three and one real one-dimensional
          block.
\end{enumerate}
The phrase ``orthogonal $2\times2$ block'' in the second case denotes the real
block for a non-real conjugate eigenvalue pair. It does not mean that the block
is an orthogonal matrix for the Euclidean scalar product.

\section{Formalized infrastructure}

The formal development contains reversal matrices, their signed variants, and
explicit congruences that expose the signatures of $P_2$, $P_3$, and $P_4$.
It also defines real Jordan blocks and realifications of non-real
conjugate-pair Jordan chains. The formalized block identities are
\begin{align}
    (\varepsilon P_n)J_n(a)
        &=J_n(a)^{\mathsf T}(\varepsilon P_n),
    \label{eq:qiclean_minkowski_canonical_pair_dependency_real_jordan_metric_identity}\\
    P_{2m}J_m(a,\tau)
        &=J_m(a,\tau)^{\mathsf T}P_{2m}.
    \label{eq:qiclean_minkowski_canonical_pair_dependency_complex_jordan_metric_identity}
\end{align}
The matrix $M$, the condition
(\ref{eq:qiclean_minkowski_canonical_pair_dependency_m_selfadjoint}), and the
four target block patterns are represented directly.\footnote{Formal
declarations in
\doclink{QICLean/Algebra/MinkowskiCanonicalPair.lean} include
\leanid{Matrix.signedSipMatrix}, \leanid{Matrix.realJordanBlock},
\leanid{Matrix.realConjugatePairJordanBlock},
\leanid{Matrix.signedSipMatrix_mul_realJordanBlock},
\leanid{Matrix.complexPairSipMatrix_mul_realConjugatePairJordanBlock},
\leanid{Matrix.minkowskiMetric}, \leanid{Matrix.IsMinkowskiSelfAdjoint}, and
\leanid{Matrix.GLRFourDimensionalPattern}.}

\section{The remaining obstruction}

The formal library has no real Jordan canonical-form existence theorem that
constructs $J$ and a similarity matrix for an arbitrary real $4\times4$
matrix. More importantly, it has no formal canonical-pair theorem that
simultaneously controls the similarity and the congruence class of the
indefinite metric, including the sign characteristic. Sylvester's law for
quadratic forms applies only after $N_J$ and the congruence in
(\ref{eq:qiclean_minkowski_canonical_pair_dependency_metric_congruence}) have
been constructed; it cannot provide the missing simultaneous reduction.

No theorem with additional hypotheses is presented as the cited result. The
existence statements
(\ref{eq:qiclean_minkowski_canonical_pair_dependency_similarity}) and
(\ref{eq:qiclean_minkowski_canonical_pair_dependency_metric_congruence}), and
the exhaustiveness of the four cases, remain unformalized. Closing the gap
requires either a source-faithful real Jordan theorem followed by a
finite-dimensional proof of the cited canonical-pair theorem, or a direct
four-dimensional proof preserving the same block, sign, transpose, and inverse
conventions. Before fixing the final formal signature, the exact existence
clause in the cited edition must also be transcribed verbatim. The present
translation of conventions is a high-confidence reconstruction from the
statement context and proof, not a verbatim quotation.

\section{Conclusion}

This note records an open formalization dependency, not a correction to either
cited source. The canonical blocks and all currently asserted algebraic
identities are proved. The simultaneous similarity--congruence existence
theorem remains the load-bearing missing result. Numerical orbit enumeration,
Euclidean singular value decomposition, or a generic indefinite-SVD assumption
would bypass rather than close the gap.

\end{document}
